\documentclass[11pt,a4paper]{article}
\usepackage[margin=25mm]{geometry}
\usepackage{lmodern}
\usepackage[T1]{fontenc}
\usepackage{amsmath,amssymb,amsthm,booktabs,url,microtype,enumitem}
\usepackage{tikz}
\usetikzlibrary{arrows.meta,positioning}
\usepackage[hidelinks]{hyperref}
\setlength{\emergencystretch}{3em}

\newtheorem{theorem}{Theorem}
\theoremstyle{remark}
\newtheorem*{remark}{Remark}

\newcommand{\lean}[1]{\texttt{#1}}
\newcommand{\cri}[1]{\texttt{#1}}
\newcommand{\uri}[1]{\texttt{#1}}
\newcommand{\rds}{\mathrm{rds}}
\newcommand{\mergef}{\mathrm{merge}}
\newcommand{\take}[1]{\mathrm{take}_{#1}}
\newcommand{\dl}{\mathrm{dropLast}}
\newcommand{\res}{\rho_{5.3}}
\newcommand{\rfc}{\rho_{3986}}
\newcommand{\conv}{\mathrm{conv}}
\newcommand{\e}{\varepsilon}
\newcommand{\seq}[1]{[#1]}
\newtheorem{lemma}{Lemma}

\title{Three Counterexamples and One-Sentence Corrections\\
for IETF Internet-Draft draft-ietf-core-href-30,\\
with a Lean Proof for the Path Rule}
\author{Nobuo Otoi\\BA Cybersecurity, Tokyo, Japan\\\texttt{otoi@ba-cybersecurity.com}}
\date{October 2026 (version 1)}

\begin{document}
\maketitle

\begin{abstract}
draft-ietf-core-href-30 defines Constrained Resource Identifiers (CRIs), a CBOR form of URIs, with a
procedure that resolves a CRI \emph{reference} against a base CRI (Section~5.3) and a procedure that
converts a CRI reference to a URI reference (Section~6.1), which the draft calls the equivalent URI
reference. For a base and a reference we compare two computations: the URI of the Section~5.3 result, and
the RFC~3986 resolution of the converted reference against the base. We give three inputs on which the two
differ, with every step written out. (A)~Resolving \cri{["t", null, ["p"]]} (the CRI for \uri{t:/p})
against \cri{s://h} yields \cri{t://h/p} by Section~5.3, whereas Section~2.1 assigns no authority to the
equivalent URI reference and RFC~3986 gives a reference with a scheme its own authority; the same input,
used as a reference to itself, also fails the draft's requirement that resolution yield the original CRI.
(B)~The reference \cri{[1]}, which removes the last path segment of the base, is converted to the empty
URI reference, which denotes the base itself. (C)~The reference \cri{[1, ["", "a"]]} is converted to
\uri{/a}, an absolute path, whereas Section~5.3 gives \uri{/b//a} for the base path \uri{/b/c}. We propose
one sentence of correction for each. On an enumeration of 720{,}120 base--reference pairs with an
authority, the two computations differ on 93{,}810; after the corrections they differ on none of the
713{,}660 pairs that remain convertible (the corrections refuse 10{,}540 pairs and admit 4{,}080 that
the draft's syntax check rejects). In Lean~4 (core library; axioms \lean{propext} and \lean{Quot.sound}
only; no \lean{sorry}) the counterexamples are checked by kernel evaluation, and the corrected conversion
of a discard-form reference, resolved per RFC~3986, is proved equal to the Section~5.3 resolution for
every base with an authority and paths of any length. We make no claim about bases without an authority.
\end{abstract}

\section{Introduction}
\label{sec:intro}

A CRI~\cite{href} carries the five components of a URI~\cite{rfc3986} as a CBOR~\cite{rfc8949} array.
A \emph{CRI reference} is the compact form for relative links: instead of the leading dot segments and
slash of a URI reference, it carries a \emph{discard} value that says how many trailing segments of the
base path to remove. Two procedures on references are normative. Section~5.3 resolves a reference
against a base CRI (``An application MUST implement this operation by applying the algorithm specified
in Section~5.3 (or any algorithm that is functionally equivalent to it)'', Section~5), and Section~6.1
converts a CRI reference to a URI reference (``Applications MUST convert a CRI reference to a URI
reference by determining the components \ldots{} according to the following steps''). The draft
tabulates the correspondence between discard values and URI references (its Table~1:
\cri{[1, ["a"]]} is \uri{a}, \cri{[2, ["a"]]} is \uri{../a}, \cri{[true, ["a"]]} is \uri{/a}).

This note compares, for a base and a reference, the URI of the Section~5.3 result with the RFC~3986
resolution of the Section~6.1 conversion against the base (Section~\ref{sec:cex} defines the comparison),
and reports three inputs on which the two differ. Each is a finite computation that the reader can
redo from the quoted text in a few minutes (Section~\ref{sec:cex}). We propose a one-sentence
correction for each (Section~\ref{sec:fix}) and support the corrections in two ways
(Section~\ref{sec:checks}): an exhaustive enumeration of small inputs, and a Lean~4 proof that the
corrected conversion agrees with RFC~3986 resolution for all path lengths. Section~\ref{sec:scope}
states exactly what is and is not claimed.

\section{The text under test}
\label{sec:bg}

\paragraph{RFC 3986.} Section~5.2.2 resolves a reference $R$ against a base $B$ (which has no fragment,
Section~5.1). If $R$ has a scheme, the target is $R$ with dot segments removed and in particular has
$R$'s authority (possibly none). Otherwise the scheme is $B$'s; if $R$ has an authority, authority and path
are $R$'s; otherwise the authority is $B$'s and the path is $B$'s path when $R$'s path is empty,
$\rds(R.\mathit{path})$ when it starts with ``/'', and $\rds(\mergef(B,R))$ otherwise, where $\mergef$
appends $R$'s path to $B$'s path with its last segment removed (or to ``/'' when $B$ has an authority
and an empty path). The query is $R$'s, except that an empty-path reference without a query keeps $B$'s
query; the fragment is $R$'s.

\paragraph{Notation.} We write a path as a list of segments $p=\seq{p_1,\dots,p_k}$, $|p|=k$, with
concatenation $p\mathbin{+\!\!+}q$, $\take{j}(p)=\seq{p_1,\dots,p_{\min(j,k)}}$ and
$\dl(p)=\take{k-1}(p)$; $\e$ is the empty segment. For a base $B$ with an authority and path $p_B$,
Sections~5.2.3 and~5.2.4 of RFC~3986 become, on segment lists,
\begin{align*}
\mergef_B(p) &= \begin{cases} p & \text{if } p_B=\seq{}\\ \dl(p_B)\mathbin{+\!\!+}p & \text{otherwise,}\end{cases}
\qquad\qquad \rds(p)=R(\seq{},p),\\
R(o,\seq{}) &= o,\\
R(o,\text{``.''}::t) &= \begin{cases} o\mathbin{+\!\!+}\seq{\e} & t=\seq{}\\ R(o,t) & t\neq\seq{},\end{cases}\\
R(o,\text{``..''}::t) &= \begin{cases} \dl(o)\mathbin{+\!\!+}\seq{\e} & t=\seq{}\\ R(\dl(o),t) & t\neq\seq{},\end{cases}\\
R(o,s::t) &= R(o\mathbin{+\!\!+}\seq{s},t)\quad\text{for any other segment } s.
\end{align*}
A URI reference without scheme or authority is $u=(\mathit{rooted},p,q,f)$ with $p$ a segment list or
$\bot$ (the empty path) and $q,f$ possibly $\bot$. Section~5.2.2 then reads
\[
\rfc(B,u).\mathit{path}=\begin{cases} p_B & p=\bot\\ \rds(p) & \mathit{rooted}\\ \rds(\mergef_B(p)) & \text{otherwise,}\end{cases}
\qquad
\rfc(B,u).\mathit{query}=\begin{cases} q & q\neq\bot\\ q_B & q=\bot,\ p=\bot\\ \seq{} & q=\bot,\ p\neq\bot,\end{cases}
\]
and the fragment is $f$. (The string form of a rooted $p$ is ``/''$p_1$``/''$\cdots p_k$; of an
unrooted one, $p_1$``/''$\cdots p_k$; $\bot$ is the empty string.) These definitions are the ones
formalized in Lean (Section~\ref{sec:checks}); their agreement with the string algorithms of RFC~3986
is checked in Section~\ref{sec:checks}.

\paragraph{The CRI data model.} A CRI reference (Figure~2 of~\cite{href}, in CDDL~\cite{rfc8610}) is an
array whose head is either a scheme and an authority or a discard value, followed by path, query and
fragment. The authority is a host, or one of two ``no authority'' values: null for a root-based path
(\uri{a:/p}) and true for a rootless path (\uri{a:p}). Table~2 of Section~5.1 gives null as the default
value of the authority section. Section~2.1 says of references:
\begin{quote}\small
(For example, for a CRI reference where authority is either not set or has either of the NOAUTHORITY
values, the equivalent URI reference's authority is absent.)
\end{quote}
Path segments are never ``.'' or ``..'' (constraint C9).

\paragraph{Section 5.3 (resolution).} Steps~3 and~5 read:
\begin{quote}\small
3. If the value of discard is true in the CRI reference (which is implicitly the case when scheme
and/or authority are present in the reference), replace the path in the buffer with the empty array,
set query to empty and fragment to null, and set a true authority to null. If the value of discard is
an unsigned integer, remove as many elements from the end of the path array; if it is non-zero, set
query to empty and fragment to null. Set discard to true in the buffer.\\[2pt]
5. If query is non-null in the CRI reference, set fragment to null in the buffer. Apart from the
path and discard, copy all non-null sections from the CRI reference to the buffer in sequence.
\end{quote}
Step~4 reads: ``If the path section is non-null in the CRI reference, append all elements from the path
array to the array in the path section in the buffer; set query to empty and fragment to null.''
For a reference in discard form $r=(d,\pi,\kappa,\varphi)$ ($d\in\{\mathit{true}\}\cup\mathbb{N}$;
$\pi$ a segment list or $\bot$; $\kappa$ a query or $\bot$; $\varphi$ a fragment or $\bot$) and a base
$B$ with path $p_B$ and query $q_B$, steps~3--5 compute
\begin{align*}
P_3&=\begin{cases}\seq{} & d=\mathit{true}\\ \take{|p_B|-n}(p_B) & d=n,\end{cases}
&Q_3&=\begin{cases} q_B & d=0\\ \seq{} & \text{otherwise,}\end{cases}\\
\res(B,r).\mathit{path}&=\begin{cases} P_3\mathbin{+\!\!+}\pi & \pi\neq\bot\\ P_3 & \pi=\bot,\end{cases}
&\res(B,r).\mathit{query}&=\begin{cases} \kappa & \kappa\neq\bot\\ \seq{} & \kappa=\bot,\ \pi\neq\bot\\ Q_3 & \kappa=\bot,\ \pi=\bot,\end{cases}
\end{align*}
and $\res(B,r).\mathit{fragment}=\varphi$.

\paragraph{Section 6.1 (conversion).} The path rule reads, in part:
\begin{quote}\small
If the CRI reference contains a discard item of value true, the path component is considered rooted.
If it contains a discard item of value 0 and the path item is present, the conversion fails. If it
contains a positive discard item, the path component is considered unrooted and prefixed by as many
``../'' components as the discard value minus one indicates. If the discard value is 1 and the first
element of the path contains a :, the path component is prefixed by ``./''. \ldots{} If the CRI
reference contains one or more path items, the path component is constructed by concatenating the
sequence of representations of these items. These representations generally contain a leading slash
(``/'') character and the value of each item, processed as discussed below. The leading slash character
is omitted for the first path item only if the path component is considered ``unrooted''.
\end{quote}
The rule ends with syntax checks: when neither authority nor scheme is present, the path component must
match path-absolute, path-noscheme or path-empty of RFC~3986, else the conversion fails.
In the notation above, with $\mathit{items}=\pi$ (or $\seq{}$ if $\pi=\bot$), the path part of
$\conv(r)$ is
\begin{align*}
d=\mathit{true}&:\ (\mathit{rooted},\ \mathit{items});\\
d=0&:\ \text{fail if }\pi\neq\bot,\ \text{else the empty path};\\
d=n\ge1&:\ (\mathit{unrooted},\ \mathit{pre}\mathbin{+\!\!+}\mathit{items}),\qquad
\mathit{pre}=\begin{cases}\seq{\text{``.''}} & n=1 \text{ and ``:''}\in\mathit{items}_1\\ \seq{\text{``..''}}^{\,n-1} & n\ge2\\ \seq{} & \text{otherwise,}\end{cases}
\end{align*}
followed by the string rendering, and the query part is $\kappa$ if $\kappa\neq\bot$ and
$\kappa\neq\seq{}$, else $\bot$.
Section~2.3 states that some well-formed references cannot be converted, naming \cri{[0, ["p"]]} and
\cri{[0, null, []]}.

\section{Three counterexamples}
\label{sec:cex}

For a base $B$ and a reference $r$ we compute two things: $u_1$, the URI of the Section~5.3 result
(Section~6.1 applied to the resolved full CRI); and $u_2$, the RFC~3986 resolution, against the URI of
$B$, of the URI reference that Section~6.1 produces from $r$---the reference that the draft calls the
equivalent URI reference (Section~2.1, Table~1). We say that $B$ and $r$ \emph{disagree} when
$u_1\neq u_2$:
\[
u_1(B,r)=\conv(\res(B,r)),\qquad u_2(B,r)=\rfc\bigl(\conv(B),\conv(r)\bigr).
\] This describes two computations; it does not presuppose that the draft promises their
agreement. In each case below we write the base and the reference, apply Section~5.3 step by step, apply
Section~6.1 step by step, and apply RFC~3986 to the converted reference. Table~\ref{tab:cex} summarizes.

\begin{table}[ht]
\centering\small
\begin{tabular}{lllllll}
\toprule
 & base & CRI reference & Section 5.3 gives & Section 6.1 gives & RFC 3986 gives\\
\midrule
A & \cri{s://h} & \cri{["t", null, ["p"]]} & \cri{t://h/p} & \uri{t:/p} & \uri{t:/p}\\
B & \cri{s://h/} & \cri{[1]} & \cri{s://h} & \uri{""} & \uri{s://h/}\\
B & \cri{s://h/a/} & \cri{[2, []]} & \cri{s://h} & \uri{../} & \uri{s://h/}\\
C & \cri{s://h/b/c} & \cri{[1, ["", "a"]]} & \cri{s://h/b//a} & \uri{/a} & \uri{s://h/a}\\
B$'$ & \cri{s://h?q} & \cri{[0, null, []]} & \cri{s://h} & \uri{""} & \uri{s://h?q}\\
\bottomrule
\end{tabular}
\caption{The counterexamples. ``Section 5.3 gives'' is the resolved CRI, written as a URI; ``Section 6.1 gives'' is the
URI reference obtained from the CRI reference; ``RFC 3986 gives'' is that URI reference resolved against the base.}
\label{tab:cex}
\end{table}

\subsection{A: a null authority is not copied}
Base \cri{["s", ["h"]]} (\uri{s://h}); reference \cri{["t", null, ["p"]]}, the CRI for \uri{t:/p}
(scheme \cri{t}, authority null = NOAUTH-ROOTBASED of Figure~2, path \cri{["p"]}).
\begin{itemize}[nosep]
\item Step~3: a scheme is present, so discard is implicitly true. The buffer's path becomes \cri{[]},
query empty, fragment null. The buffer's authority is \cri{["h"]}, not true, so it is unchanged.
\item Step~4: the path \cri{["p"]} is appended: buffer path \cri{["p"]}.
\item Step~5: the non-null sections of the reference are the scheme \cri{t} and the path. The authority
is null and is not copied. Result: \cri{["t", ["h"], ["p"]]}, i.e.\ \uri{t://h/p}.
\end{itemize}
In symbols, with $B=(\mathit{s},\seq{h},\seq{},\seq{})$ and $r=(\mathit{t},\mathit{null},\seq{p})$:
\[
\res(B,r)=(\mathit{t},\seq{h},\seq{p}),\qquad\text{RFC 3986 for a reference with a scheme: } T=R=(\mathit{t},\bot,\seq{p}).
\]
Section~2.1 assigns no authority to the equivalent URI reference of a reference whose authority is null
(or not set), and RFC~3986 Section~5.2.2 gives a reference with a scheme its own authority: \uri{t:/p}.
The Section~5.3 result differs from both. The outcome does not depend on how null is read (a value, as
the CDDL has it, or ``not set''), since Section~2.1 covers both; nor on whether ``a true authority'' in
step~3 refers to the buffer's or the reference's authority, since neither is true.

A second check stays inside the draft. Section~5 says that ``the only requirement on a CRI reference is
that reference resolution MUST yield the original CRI'', and Figure~2 admits a reference with a scheme and
a null authority, so a CRI may be written as a reference to itself. Applying steps~3--5 to the CRI
\uri{t:/p}, used as a reference against the base \uri{s://h}, yields \uri{t://h/p}, not the original;
the rootless \uri{t:p} (authority true), by contrast, is returned unchanged, because step~5 copies true.

\subsection{B: a discard value without path items}
Base \cri{["s", ["h"], [""]]} (\uri{s://h/}); reference \cri{[1]}.
\begin{itemize}[nosep]
\item Section~5.3: discard~1 removes the last element of the path, \cri{[""]} becomes \cri{[]}; query
empty, fragment null; nothing to append or copy. Result \cri{["s", ["h"], []]}, i.e.\ \uri{s://h}.
\item Section~6.1: discard is positive, so the path is unrooted and prefixed by $1-1=0$ copies of
\uri{../}. The rule then defines the path component only ``if the CRI reference contains one or more
path items''; here there are none. Read as an empty path component (the reading of our model), the URI
reference is \uri{""}. Read as undefined, Section~6.1 produces nothing and, unlike the case of discard~0
with a path, does not say that the conversion fails.
\item RFC~3986, under the first reading: the empty reference has an empty path, so the target path is the
base path: \uri{s://h/}.
\end{itemize}
In symbols, with $p_B=\seq{\e}$ and $r=(1,\bot,\bot,\bot)$:
\begin{gather*}
\res(B,r).\mathit{path}=\take{0}(\seq{\e})=\seq{},\qquad
\conv(r)=(\mathit{unrooted},\seq{})\ \text{rendered as }\bot,\\
\rfc(B,\conv(r)).\mathit{path}=p_B=\seq{\e}.
\end{gather*}
Under either reading the failure clause of correction~(B) is missing. The same happens for \cri{[1, []]}
and \cri{[true]} (converted to \uri{""}), and for \cri{[2, []]} against \uri{s://h/a/}: Section~5.3
removes both segments and gives \uri{s://h}, Section~6.1 gives \uri{../}, and RFC~3986 resolves
\uri{../} against \uri{s://h/a/} to \uri{s://h/}:
\begin{align*}
\res(B,r).\mathit{path}&=\take{0}(\seq{a,\e})\mathbin{+\!\!+}\seq{}=\seq{},\\
\rds(\mergef_B(\seq{\text{``..''}}))&=R(\seq{},\seq{a,\text{``..''}})=R(\seq{a},\seq{\text{``..''}})\\
&=\dl(\seq{a})\mathbin{+\!\!+}\seq{\e}=\seq{\e}.
\end{align*}
A relative-path reference cannot have the meaning of \cri{[1]} for all bases:
\begin{lemma}
For every base $B$ with an authority and every non-empty segment list $p$,
$\mathrm{last}\bigl(\rds(\mergef_B(p))\bigr)=\e$ if $\mathrm{last}(p)\in\{\e,\text{``.''},\text{``..''}\}$ and
$=\mathrm{last}(p)$ otherwise; and $\rfc(B,(\cdot,\bot,\bot,\bot)).\mathit{path}=p_B$.
\end{lemma}
\noindent (In $R$, the last segment of the input is the last one appended, or $\e$ in the three
special cases; nothing after it is removed.) The last segment of the target therefore does not depend
on $B$, while a reference meaning ``remove the last segment'' would have to end in $B$'s second-to-last
segment. Section~2.3 states that some references
cannot be converted, and Section~6.1 fails for \cri{[0, ["p"]]}; it does not fail for these.

The reference \cri{[0, null, []]} (B$'$ in Table~\ref{tab:cex}; ``leave the path, remove the query''),
which Section~2.3 lists among the references that cannot be converted, is likewise converted to
\uri{""}, which keeps the base's query under RFC~3986.

\subsection{C: an empty first segment}
Base \cri{["s", ["h"], ["b", "c"]]} (\uri{s://h/b/c}); reference \cri{[1, ["", "a"]]}.
\begin{itemize}[nosep]
\item Section~5.3: discard~1 removes \cri{"c"}; the items \cri{""} and \cri{"a"} are appended. Result
\cri{["s", ["h"], ["b", "", "a"]]}, i.e.\ \uri{s://h/b//a}.
\item Section~6.1: discard~1, unrooted, no \uri{../} prefix; the first element does not contain a colon,
so no \uri{./} prefix. The first item \cri{""} is written without a leading slash, the second as
\uri{/a}: the path component is \uri{""}\,+\,\uri{/a} = \uri{/a}. It matches path-absolute, so the
conversion succeeds with the URI reference \uri{/a}.
\item RFC~3986: a reference whose path starts with ``/'' replaces the base path: \uri{s://h/a}.
\end{itemize}
In symbols, with $p_B=\seq{b,c}$ and $r=(1,\seq{\e,a},\bot,\bot)$:
\begin{gather*}
\res(B,r).\mathit{path}=\take{1}(\seq{b,c})\mathbin{+\!\!+}\seq{\e,a}=\seq{b,\e,a},\\
\conv(r)=(\mathit{unrooted},\seq{\e,a})\ \text{rendered as }(\mathit{rooted},\seq{a}),\qquad
\rfc(B,\conv(r)).\mathit{path}=\rds(\seq{a})=\seq{a}.
\end{gather*}
With correction~(C), $\conv(r)=(\mathit{unrooted},\seq{\text{``.''},\e,a})$ and
\[
\rds(\mergef_B(\seq{\text{``.''},\e,a}))=R(\seq{},\seq{b,\text{``.''},\e,a})=R(\seq{b},\seq{\text{``.''},\e,a})=R(\seq{b},\seq{\e,a})=\seq{b,\e,a}.
\]
A URI reference with the meaning of the CRI reference is \uri{.//a}: RFC~3986 merges it to
\uri{/b/.//a} and $\rds$ gives \uri{/b//a}. The \uri{./} prefix that Section~6.1 prescribes for a first
segment containing a colon (so that it is not read as a scheme) is also needed for an empty first segment
(so that it is not read as a leading slash). The subcase \cri{[1, [""]]}, meaning \uri{./}, is converted
to \uri{""}.

\section{Corrections}
\label{sec:fix}

\begin{enumerate}[label=\textbf{(\Alph*)},leftmargin=*]
\item Section~5.3, step~3 or step~5: ``If the CRI reference has a scheme, replace the authority in the
buffer with the authority of the CRI reference (which may be null or true).'' Equivalently, as in
RFC~3986 Section~5.2.2: ``If the CRI reference has a scheme, the result is the CRI reference itself and
steps~3--5 are skipped.''
\item Section~6.1, path rule: ``If the CRI reference contains a discard item (true or a positive
integer) and no path items, the conversion fails.'' And, for the case of Section~2.3: ``If it contains a
discard item of value~0, no path item, and a query item that is an empty array, the conversion fails.''
\item Section~6.1, path rule: ``If the discard value is~1 and the first element of the path is the empty
string or contains a :, the path component is prefixed by `./'.''
\end{enumerate}

\section{Checks}
\label{sec:checks}

\subsection{Exhaustive enumeration}
We implemented RFC~3986 Section~5.2 literally on strings (\lean{model/rfc3986.py}); it reproduces the
42 examples of Section~5.4. We implemented Sections~5.3 and~6.1 of~\cite{href} as written
(\lean{model/cri.py}), with the corrections (A)--(C) behind a flag. For every base CRI with a host
authority and every CRI reference in an enumeration (path segments from \cri{"a"}, \cri{"b"}, the empty
string and \cri{"x:y"}; base paths of up to three segments; references with discard~0--4 or true, or
with a scheme or an authority; queries and fragments present or absent), we compared the URI of the
Section~5.3 result with the RFC~3986 resolution of the Section~6.1 conversion, on the pairs that
Section~6.1 converts (\lean{model/check.py}). Where a discard value exceeds the number of base path
segments, we read ``remove as many elements'' as stopping at the empty array (the text does not say).
With $\mathcal{B}$ the 170 bases and $\mathcal{R}$ the 4{,}890 references, let
$\mathcal{C}=\{(B,r)\in\mathcal{B}\times\mathcal{R}:\conv(r)\text{ defined}\}$ and
$\mathcal{D}=\{(B,r)\in\mathcal{C}:u_1(B,r)\neq u_2(B,r)\}$; the corrected versions are
$\mathcal{C}'$ and $\mathcal{D}'$. Then $|\mathcal{B}\times\mathcal{R}|=831{,}300$,
$|\mathcal{C}|=720{,}120$, $|\mathcal{D}|=93{,}810$, $|\mathcal{C}'|=713{,}660$ and $|\mathcal{D}'|=0$,
with the classification of $\mathcal{D}$ by the shape of the reference (B comprises the ten classes
discard 1--4 or true with the path absent or \cri{[]}):

\begin{center}\small
\begin{tabular}{llrr}
\toprule
 & & draft -30 & corrected\\
\midrule
\multicolumn{2}{l}{comparable pairs} & 720{,}120 & 713{,}660\\
A & scheme present, authority null & 66{,}300 & 0\\
B & positive or true discard, no path items & 10{,}168 & 0\\
B$'$ & discard 0, no path, query \cri{[]} & 170 & 0\\
C & discard 1, first segment empty & 17{,}172 & 0\\
\multicolumn{2}{l}{total disagreements} & 93{,}810 & 0\\
\bottomrule
\end{tabular}
\end{center}
\noindent Every disagreement is of one of the three shapes, and none remains after the corrections.
The number of comparable pairs changes because correction~(B) refuses 10{,}540 of the 720{,}120 pairs,
while the corrected rule~(C) admits 4{,}080 pairs of the shape \cri{[1, ["", "", \ldots]]} (converted
to \uri{.//\ldots}) that the draft's syntax check rejects (\uri{//\ldots}). A smaller enumeration (paths of up to two segments, discard up to~3; 40{,}320 comparable pairs, 7{,}530
disagreements) gives the same picture.

\subsection{Lean 4}
\lean{lean/Cri.lean} (Lean~4.33.1~\cite{lean4}, core library only, 328 lines) models paths as lists of
segments, where a segment is abstracted to what the procedures inspect: \lean{Seg = dot | dotdot | empty
| other id colon} (an identifier and whether the segment contains a colon). \lean{resolveCRI} is
steps~3--5 of Section~5.3 for references in discard form; \lean{convDraft} is the Section~6.1 path rule
under the empty-path reading of Section~\ref{sec:cex}, followed by the collapse that the string form
imposes on an unrooted path whose first segment is empty (\uri{""}\,+\,\uri{/a} is \uri{/a}, \uri{""}
alone is the empty path); it omits the final path-syntax check of Section~6.1, which none of the four
counterexamples triggers (the theorem concerns \lean{convFixed}). \lean{convFixed}
applies (B) and (C); \lean{rds}, \lean{merge} and \lean{resolve3986} are RFC~3986 Sections~5.2.2--5.2.4
for a base with an authority and a reference without scheme or authority.

\paragraph{Counterexamples by kernel evaluation.} Four theorems each state that \lean{convDraft}
yields the URI reference of Table~\ref{tab:cex} and that its RFC~3986 resolution differs from
\lean{resolveCRI}:
\begin{itemize}[nosep]
\item \lean{cex\_discard\_without\_path}: \cri{[1]} against \uri{s://h/} (B);
\item \lean{cex\_discard\_two\_without\_path}: \cri{[2, []]} against \uri{s://h/a/} (B);
\item \lean{cex\_leading\_empty\_segment}: \cri{[1, ["", "a"]]} against \uri{s://h/b/c} (C);
\item \lean{cex\_discard\_zero\_empty\_query}: \cri{[0, null, []]} against \uri{s://h?q} (B$'$).
\end{itemize}
All are proved by \lean{decide} and depend on \lean{propext} only.
Finding~A concerns the authority, which the path model abstracts away; \lean{authDraft} models steps~3
and~5 on the three authority values, and \lean{cex\_null\_authority} states
\lean{authDraft (host 0) rootbased = host 0} with no axioms.

\begin{theorem}[\lean{convFixed\_sound}]
Let $B$ be a base with an authority whose path contains no dot segments, and $r=(d,\pi,\kappa,\varphi)$ a
CRI reference in discard form with $\pi$ free of dot segments. Then
\[
\conv'(r)=u\ \Longrightarrow\ \rfc(B,u)=\res(B,r),
\]
where $\conv'$ is the conversion with corrections (B) and (C) (\lean{convFixed}).
\end{theorem}
\begin{proof}[Proof sketch]
By cases on the discard value. For true, the reference has a non-empty path $p$ and $\rds(p)=p$. For~0,
the reference has no path and its query is not \cri{[]}; both sides keep the base path and take the
reference's query if present, else the base's. For $k+1$, the reference has a non-empty path $s{::}t$ and
the converted path is $\mathit{pre}\mathbin{+\!\!+}(s{::}t)$ with $\mathit{pre}$ empty or $[\text{``.''}]$
for $k=0$ and $k$ copies of ``..'' otherwise. The lemmas \lean{rdsAux\_dotdots} and \lean{rdsAux\_dot} give, for non-empty dot-free
$\mathit{items}$ and any $o$,
\begin{gather*}
R\bigl(o,\ \seq{\text{``..''}}^{k}\mathbin{+\!\!+}\mathit{items}\bigr)=\take{|o|-k}(o)\mathbin{+\!\!+}\mathit{items},\qquad
R\bigl(o,\ \text{``.''}::\mathit{items}\bigr)=o\mathbin{+\!\!+}\mathit{items}
\end{gather*}
(the first by induction on $k$, using $\dl(o)=\take{|o|-1}(o)$).
Since $\mergef_B$ removes one more segment,
\[
\rds\bigl(\mergef_B(\mathit{pre}\mathbin{+\!\!+}(s{::}t))\bigr)=\take{|p_B|-(k+1)}(p_B)\mathbin{+\!\!+}(s{::}t)=P_3\mathbin{+\!\!+}\pi ,
\]
which is step~3 followed by step~4 (the case $p_B=\seq{}$, where $\mergef_B$ is the identity and
$\take{|p_B|-(k+1)}$ is empty, is immediate). Query and fragment agree by inspection of the two
definitions.
\end{proof}
\noindent The theorem depends on \lean{propext} and \lean{Quot.sound}; \lean{Classical.choice} is not
used and the file contains no \lean{sorry}. \lean{lean Cri.lean} prints only the six
\lean{\#print axioms} lines. The theorem covers references in discard form; correction~(A) concerns the
authority, which this model abstracts away, and rests on the enumeration.

\paragraph{Relation between the two checks.} The segment-level rendering of RFC~3986 in Lean was
compared with the literal string implementation on all 1{,}364 rooted paths of one to five segments over
\{\cri{a}, empty, \cri{.}, \cri{..}\} and on all 21{,}420 (draft) and 19{,}320 (corrected) convertible
pairs of a host base and a discard-form reference of the smaller enumeration, with no difference
(\lean{model/seglevel.py}); \lean{lean/CrossCheck.lean} re-checks by \lean{decide} the results that the
Python model computed for 60 random pairs.

\section{A remark on host names}
\label{sec:ipv4}
Section~3 of~\cite{href} states the objective that two distinct normalized CRIs do not convert to the
same URI. Constraint~C5 allows the host name \cri{["1","2","1","2"]}, which converts to \uri{1.2.1.2},
as does the IP-address host \cri{h'01020102'}; RFC~3986 Section~3.2.2 says that a host matching
IPv4address is an IPv4 address literal, and Section~8.1.2 of~\cite{href} maps such a Uri-Host to
host-ip. A sentence in C5 (``A registered name that matches the IPv4address rule MUST be represented as
host-ip'') would make the data model say the same.

\section{Scope}
\label{sec:scope}
We claim exactly the following. (1)~On the inputs of Table~\ref{tab:cex}, Sections~5.3 and~6.1 of
draft-ietf-core-href-30, applied as written, give the results shown, which differ from the authority that
Section~2.1 and RFC~3986 give (A) or from the RFC~3986 meaning of the converted reference (B, C, B$'$).
(2)~With corrections (A)--(C), no such disagreement exists on the enumeration of
Section~\ref{sec:checks}, and the Lean theorem extends (B) and (C) to all path lengths for bases with an
authority and discard-form references, at the segment level.
We do not claim anything about base CRIs without an authority (there RFC~3986's own $\rds$ turns
rootless paths into rooted ones, so we leave them aside), about the
practical frequency of the shapes involved, or about consequences beyond the plain fact that a
converted reference denotes a different URI than the CRI reference. The segment-level model is related
to the string procedures by testing, not by proof. The reference shape \cri{[null, true, path]} was
excluded from the enumeration.

In the 144 closed issues and pull requests of \lean{core-wg/href} and the messages of the \lean{core}
mailing list that we examined (listed in \lean{prior\_art.md}) (2020 onwards, about 90 matching \emph{href}, \emph{CRI},
\emph{remove\_dot\_segments}, \emph{rootless} or \emph{discard}; twelve read in full), none of A, B or C
is discussed; issue~\#48 introduced the \uri{./} prefix for a leading colon, and issue~\#30 found the
wording of step~5 odd without drawing consequence~A.

\section*{Artifacts}
\lean{lean/Cri.lean}, \lean{lean/CrossCheck.lean} (Lean~4.33.1, core only) and \lean{model/}
(Python~3: \lean{rfc3986.py}, \lean{cri.py}, \lean{check.py}, \lean{inject.py}, \lean{seglevel.py},
\lean{gen\_crosscheck.py}, with their outputs, produced by re-running the scripts in the archive), and
\lean{prior\_art.md} (the issues and messages examined). Checked against the text of
draft-ietf-core-href-30 (21 November 2025). Repository and archive: to be published with this note.

\section*{Acknowledgements}
An AI agent, applying a target-selection method set by the author, chose this draft and produced the
proofs, code and text. The author restricted the claims to those with self-evident evidence, directed
two rounds of revision, and checked and approved the results.

\begin{thebibliography}{9}\small
\bibitem{href} C.~Bormann (Ed.) and H.~Birkholz, ``Constrained Resource Identifiers,''
  Internet-Draft draft-ietf-core-href-30, IETF, 21 November 2025; in the RFC Editor queue as of
  October 2026.
\bibitem{rfc3986} T.~Berners-Lee, R.~Fielding, and L.~Masinter, ``Uniform Resource Identifier (URI):
  Generic Syntax,'' STD~66, RFC~3986, January 2005.
\bibitem{rfc8949} C.~Bormann and P.~Hoffman, ``Concise Binary Object Representation (CBOR),'' STD~94,
  RFC~8949, December 2020.
\bibitem{rfc8610} H.~Birkholz, C.~Vigano, and C.~Bormann, ``Concise Data Definition Language (CDDL),''
  RFC~8610, June 2019.
\bibitem{lean4} L.~de~Moura and S.~Ullrich, ``The Lean 4 Theorem Prover and Programming Language,''
  in \emph{CADE 28}, LNCS 12699, Springer, 2021.
\end{thebibliography}

\end{document}
